\section*{Chapter \ref{ch:g8:ice-on-earth} --- Ice on Earth}

\begin{solution}{exo:g8:ice-on-earth:1}
(a) The accumulation area. (b) The ablation area. (c) The snout. (d) A
(terminal) \omterm{def:g8:ice-on-earth:glacial-landforms}{moraine}.
\end{solution}

\begin{solution}{exo:g8:ice-on-earth:2}
(a) False: a \omterm{def:g8:ice-on-earth:glacier}{glacier} is made of \omterm{def:g4:clouds-rain-snow:precipitation}{snow} fallen on land and turned into ice;
frozen \omterm{def:g2:water-on-earth:ocean}{sea} water is \omterm{def:g8:ice-on-earth:sea-ice}{sea ice}. (b) True. (c) False: its ice still flows
downhill; the snout moves back because the ice melts faster than it
flows. (d) True.
\end{solution}

\begin{solution}{exo:g8:ice-on-earth:3}
(a) A \omterm{def:g8:ice-on-earth:glacial-landforms}{U-shaped valley}. (b) An erratic. (c) An \omterm{def:g8:ice-on-earth:ice-sheet}{ice shelf}. (d) \omterm{def:g8:ice-on-earth:permafrost}{Permafrost}.
\end{solution}

\begin{solution}{exo:g8:ice-on-earth:4}
(a) Antarctica and Greenland. (b) Antarctica, by about \qty{58}{m}
(Greenland: about \qty{7.4}{m}). (c) About a tenth (10~\%).
\end{solution}

\begin{solution}{exo:g8:ice-on-earth:5}
(a) About \qty{7.05}{million~km^2} in 1979 and \qty{3.57}{million~km^2}
in 2012. (b) $7.05 - 3.57 = 3.48$: about \qty{3.5}{million~km^2}.
(c) $3.48 \div 7.05 \approx 0.49$: about half. (d) A single year depends
on the \omterm{def:g1:watching-the-weather:weather}{weather} of that summer: 2012 is far below the line. The trend line
uses all the years, so one unusual summer changes it very little.
\end{solution}

\begin{solution}{exo:g8:ice-on-earth:6}
(a) $2004 - 1941 = 63$ years. (b) $12000 \div 63 \approx 190.5$: about
\qty{190}{m} a year. (c) $2021 - 1984 = 37$ years; $3200 \div 37 \approx
86$: about \qty{86}{m} a year. (d) Muir \omterm{def:g8:ice-on-earth:glacier}{Glacier}, more than twice as fast.
\end{solution}

\begin{solution}{exo:g8:ice-on-earth:7}
In the \emph{accumulation area}, high on the \omterm{def:g1:landscapes-around-us:landform}{mountain}, more \omterm{def:g4:clouds-rain-snow:precipitation}{snow} falls
each year than melts. New \omterm{def:g4:clouds-rain-snow:precipitation}{snow} buries and presses the old \omterm{def:g4:clouds-rain-snow:precipitation}{snow}, which
turns into \emph{firn} after about a year, then, deeper down and much
later, into solid ice. Under its own weight the ice \emph{flows} slowly
down the \omterm{def:g1:landscapes-around-us:landform}{valley}. Lower down, in the \emph{ablation area}, it melts faster
than \omterm{def:g4:clouds-rain-snow:precipitation}{snow} falls; the ice that reaches the snout is ice that formed high
up, often long ago.
\end{solution}

\begin{solution}{exo:g8:ice-on-earth:8}
(a) In glass B: the melted ice runs into the water and raises it; in
glass A the level stays where it was. (b) Glass A is the \omterm{def:g8:ice-on-earth:ice-sheet}{iceberg} (ice
already floating in the \omterm{def:g2:water-on-earth:ocean}{sea}); glass B is the \omterm{def:g8:ice-on-earth:ice-sheet}{ice sheet} (ice on land).
(c) An \omterm{def:g8:ice-on-earth:ice-sheet}{ice shelf} holds back the \omterm{def:g8:ice-on-earth:glacier}{glaciers} behind it. When it breaks up,
these \omterm{def:g8:ice-on-earth:glacier}{glaciers} speed up and pour more land ice into the \omterm{def:g2:water-on-earth:ocean}{sea}, and that
raises it.
\end{solution}

\begin{solution}{exo:g8:ice-on-earth:9}
(a) A \omterm{def:g8:ice-on-earth:glacial-landforms}{U-shaped valley}, carved by a \omterm{def:g8:ice-on-earth:glacier}{glacier}. (b) A \omterm{def:g8:ice-on-earth:glacier}{glacier} carried it from
far up the \omterm{def:g1:landscapes-around-us:landform}{valley}, where there is granite, and left it when the ice
melted: it is an erratic. (c) A lateral \omterm{def:g8:ice-on-earth:glacial-landforms}{moraine}: its height shows how
thick the \omterm{def:g8:ice-on-earth:glacier}{glacier} once was; the ice reached at least that high on the
\omterm{def:g1:landscapes-around-us:landform}{valley} side.
\end{solution}

\begin{solution}{exo:g8:ice-on-earth:10}
(a) $7.4 \div 10 = \qty{0.74}{m}$. (b) $58 \div 100 = \qty{0.58}{m}$.
(c) $58 \div 7.4 \approx 7.8$: about 8 times more.
\end{solution}

\begin{solution}{exo:g8:ice-on-earth:11}
(a) $7.21 \times 10^6 = \qty{7210000}{km^2}$ and $4.65 \times 10^6 =
\qty{4650000}{km^2}$. (b) $7.21 - 4.65 = \qty{2.56}{million~km^2}$, that
is \qty{2560000}{km^2}. (c) $2.56 \div 3 \approx 0.85$: about
\qty{0.85}{million~km^2} per decade, a little more than the trend line's
\qty{0.76}{million~km^2}; each decade mean still carries the ups and downs
of its own ten summers, but both ways of reading the graph show the same
large loss. (d) The extent cannot fall below zero, and a trend line only
describes the years that were measured: what happens next depends on how
the Arctic's \omterm{def:g5:weather-and-climate:climate}{climate} changes, so the line is not a forecast.
\end{solution}

\begin{solution}{exo:g8:ice-on-earth:12}
(a) $120 \div 1 = 120$: the \omterm{def:g1:landscapes-around-us:water}{coast} was about \qty{120}{km} further out to
\omterm{def:g2:water-on-earth:ocean}{sea}. (b) It had evaporated from the \omterm{def:g2:water-on-earth:ocean}{ocean}, fallen as \omterm{def:g4:clouds-rain-snow:precipitation}{snow} on the
\omterm{def:g3:continents-and-oceans:continent}{continents} and stayed there as the ice of the great \omterm{def:g8:ice-on-earth:ice-sheet}{ice sheets}. (c) A
terminal \omterm{def:g8:ice-on-earth:glacial-landforms}{moraine} forms at the end of a \omterm{def:g8:ice-on-earth:glacier}{glacier}. One on a \omterm{def:g1:landscapes-around-us:landform}{plain} far from
any \omterm{def:g1:landscapes-around-us:landform}{mountains} shows that an \omterm{def:g8:ice-on-earth:ice-sheet}{ice sheet} once reached that \omterm{def:g1:landscapes-around-us:landform}{plain}, at a time
when ice covered much more land than today.
\end{solution}

\begin{solution}{pb:g8:ice-on-earth:1}
\textbf{1.} $1080 - 0 = \qty{1080}{m}$.

\textbf{2.} 1960--1970: \qty{180}{m}; 1970--1980: \qty{150}{m};
1980--1990: \qty{190}{m}; 1990--2000: \qty{180}{m}; 2000--2010:
\qty{180}{m}; 2010--2020: \qty{200}{m}. The least in 1970--1980, the most
in 2010--2020.

\textbf{3.} No: a retreating \omterm{def:g8:ice-on-earth:glacier}{glacier} still flows downhill. Its snout moves
back because the ice melts faster than it arrives.

\textbf{4.} A terminal \omterm{def:g8:ice-on-earth:glacial-landforms}{moraine}: the ice carried \omterm{def:g3:stones-pebbles-sand:rock}{rock} debris down the
\omterm{def:g1:landscapes-around-us:landform}{valley} and dropped it at its snout, where the ice melted, while the snout
stood there.

\textbf{5.} $1080 \div 6 = \qty{180}{m}$ per decade.

\textbf{6.} $180 \div 10 = \qty{18}{m}$ a year.

\textbf{7.} $190 \div 18 \approx 10.6$: about 11 times faster.

\textbf{8.} Less: it lost more ice in its ablation area than it gained in
its accumulation area, so its snout moved back.

\textbf{9.} U-shaped: a broad floor with steep sides, because the \omterm{def:g8:ice-on-earth:glacier}{glacier}
widened and deepened the \omterm{def:g1:landscapes-around-us:landform}{valley} while it filled it.

\textbf{10.} $1980 - 1080 = \qty{900}{m}$.

\textbf{11.} $900 \div 18 = 50$ years: in \textbf{2070}.

\textbf{12.} The speed changed from decade to decade (from 150 to
\qty{200}{m}), and the future retreat depends on the snowfall and the
summer temperatures of the years to come: a past rate is not a forecast.
\end{solution}
